% ===== 原文 pp.34--39（例 3.9 方框跨 p.38--39，按 STYLE_GUIDE §7 由开始者合并写出，其收尾四式 (3.75)--(3.78) 属 p.39）=====
% [原文 p.34]
\begin{examplebox}{例 3.5：带基础表示物质的 $\SU(2)$ Higgs 化到 $\U(1)$}
给例 3.2 中的 UV 理论再添加一个新的物质场 $\psi$，它在 $\SU(2)$ 规范群的基础表示（自旋-$\frac{1}{2}$ 表示）中变换。物质场 $\psi$ 在 $\SU(2)$ 的 $\U(1)$ 极大环面下的电荷为
\begin{equation}
  q_\psi=1 \tag{3.42}
\end{equation}
因此 IR 的 1-形式对称性不再存在
\begin{equation}
  G^{(1)}=\widehat{\Z}_{\gcd(q_\Phi,q_\psi)}=0\,. \tag{3.43}
\end{equation}
这与 UV 的 1-形式对称性一致：后者（见 (3.151) 附近的讨论）在加入 $\psi$ 之后同样是平凡的。
\end{examplebox}

% [Tutor Note: p.34 式 (3.43) 原文右端写作 $=0$，而前文 (3.27) 把平凡群写作 $G^{(1)}=1$；两种都是“平凡群”的常见写法，本译本照抄不改。核对：$\gcd(q_\Phi,q_\psi)=\gcd(2,1)=1$，故 $\Z_{\gcd}=\Z_1$ 平凡、其对偶亦平凡，与“IR 的 1-形式对称性不再存在”一句自洽。]

\subsubsection{多个 阿贝尔规范场}

我们可以把上面关于电 1-形式对称性的计算推广到具有 $r$ 个 阿贝尔规范场的 阿贝尔规范理论，其规范群为
\begin{equation}
  \cG=\prod_{a=1}^{r}U(1)_a=U(1)^r\,. \tag{3.44}
\end{equation}
物质场为 $\phi_i$（$1\le i\le m$），它们的规范电荷 $q_{i,a}$ 可以组织成一个 $m\times r$ 的荷矩阵 $Q$。

Wilson 线算符取遍规范群 $\U(1)^r$ 的所有可能电荷，因而用 阿贝尔群 $\Z^r$ 的元素来标记。一个物质场 $\phi_i$ 的各电荷指定了一个向量 $Q_i\in\Z^r$。让我们通过对所有 $Q_i$ 取整系数线性组合，定义 $\Z^r$ 的一个子群 $Q_{\mathrm{matter}}$
\begin{equation}
  Q_{\mathrm{matter}}:=\mathrm{Span}_{\Z}\{Q_1,Q_2,\cdots,Q_m\}\,. \tag{3.45}
\end{equation}
于是，未屏蔽 Wilson 线的 阿贝尔群可由 $\Z^r$ 模掉 $Q_{\mathrm{matter}}$ 得到
\begin{equation}
  \mathsf{D}_1=\frac{\Z^r}{Q_{\mathrm{matter}}}\,, \tag{3.46}
\end{equation}
其 Pontryagin 对偶就是电 1-形式对称性群
\begin{equation}
  G^{(1)}=\widehat{\mathsf{D}}_1\,. \tag{3.47}
\end{equation}
计算 (3.46) 右端的一个实用办法，是利用荷矩阵 $Q$ 的 \textbf{Smith 标准形（SNF）}，它内部含有一个对角矩阵，其余元素全为零：
\begin{equation}
  \diag(q_1^D,q_2^D,\cdots,q_{\min(m,r)}^D)\subset \mathrm{SNF}(Q)\,, \tag{3.48}
\end{equation}
其中每个 $q_i^D$ 都是非负整数。

% [原文 p.35]
\begin{notebox}{Smith 标准形}
为计算一个整数矩阵 $M$ 的 Smith 标准形，我们执行下列操作的一个序列\footnote{在 Mathematica 中相应的命令是 \texttt{SmithDecomposition}。}：
\begin{itemize}[leftmargin=1.4em,itemsep=1pt,topsep=1pt]
\item 把某一行 $R_i$ 的整数倍加到另一行 $R_j$ 上，
\item 把某一列 $C_i$ 的整数倍加到另一列 $C_j$ 上，
\item 把某一行乘以负号，
\item 把某一列乘以负号，
\item 交换两行，
\item 交换两列，
\end{itemize}
直到我们到达一个整数矩阵 $\mathrm{SNF}(M)$，其中每一行、每一列都至多只有一个非零元，且任何非零元都是正的。
\end{notebox}

于是我们有
\begin{equation}
  \mathsf{D}_1=\prod_{i=1}^{\min(m,r)}\Z_{q_i^D}\,, \tag{3.49}
\end{equation}
其中在上式中 $\Z_0:=\Z$。它的 Pontryagin 对偶，也就是 1-形式对称性群，为
\begin{equation}
  G^{(1)}=\prod_{i=1}^{\min(m,r)}\widehat{\Z}_{q_i^D}\,, \tag{3.50}
\end{equation}
其中我们有
\begin{align}
  \widehat{\Z}_{q_i^D}&\cong\Z_{q_i^D}\,,\qquad q_i^D>0\,,\nonumber\\
  \widehat{\Z}_0&\cong U(1)\,. \tag{3.51}
\end{align}
另一方面，磁 $(d-3)$-形式对称性照例不受物质内容的影响：
\begin{equation}
  G^{(d-3)}=U(1)^r\,. \tag{3.52}
\end{equation}

% [Tutor Note: p.35 式 (3.49)--(3.51) 的约定 $\Z_0:=\Z$ 是整套 SNF 记账的关键：对角元为 $0$ 的那些因子对应未破缺的连续 $\U(1)$ 因子（(3.51) 第二行），而对角元 $q_i^D>0$ 给出 $\Z_{q_i^D}$。核对方法：用 (3.55) 的 $Q$ 算 SNF 得 $(3.56)$ 的对角 $(3,1)$，代入 (3.49) 得 $\Z_3\times\Z_1=\Z_3$，与 (3.57) 一致；若漏掉 $\Z_0:=\Z$ 的约定，(3.52) 的 $\U(1)^r$ 就无法与 (3.46) 的商群记账对上。]

\begin{examplebox}{例 3.6：把 $\SU(3)$ Higgs 化到 $\U(1)\times \U(1)$}
稍微推广例 3.2：考虑一个 UV 理论，其规范理论为 $\SU(3)$，并带有一个在伴随表示中的标量场 $\Phi$。取一个一般的 vev $\langle\Phi\rangle$ 会导向一个 IR 理论，其规范群是 $\SU(3)$ 的极大环面
\begin{equation}
  U(1)_1\times U(1)_2\subset SU(3)\,. \tag{3.53}
\end{equation}
% [原文 p.36]
W 玻色子的电荷为
\begin{equation}
  q_1=(2,-1)\,,\qquad q_2=(-1,2) \tag{3.54}
\end{equation}
这里的电荷是在 $\U(1)_1\times \U(1)_2$ 下测量的。于是荷矩阵为
\begin{equation}
  Q=\begin{pmatrix} 2 & -1\\ -1 & 2\end{pmatrix}\,. \tag{3.55}
\end{equation}
它的 Smith 标准形为
\begin{equation}
  \mathrm{SNF}(Q)=\begin{pmatrix} 3 & 0\\ 0 & 1\end{pmatrix}\,. \tag{3.56}
\end{equation}
由此我们算出 IR 的 1-形式对称性群为
\begin{equation}
  G^{(1)}=\widehat{\Z}_3\times\widehat{\Z}_1=\Z_3\times\Z_1=\Z_3\,, \tag{3.57}
\end{equation}
这与后面将要讨论的 UV 1-形式对称性群相匹配（见 (3.153) 附近的讨论）。
\end{examplebox}

% [Tutor Note: p.35--36 例 3.6 可用来验证 SNF 的“整系数换基不改变商群”这一事实：$Q$ 的两个 W 玻色子电荷 $(2,-1),(-1,2)$ 是 $\Z^2$ 中的格点，(3.55) 的行列式为 $3$，故 $\Z^2/\mathrm{Span}_\Z\{Q_i\}$ 的阶为 $3$，即 $\Z_3$，与 (3.57) 一致。注意 $\widehat{\Z}_1=\Z_1$（平凡群），所以 (3.57) 中间一步的 $\Z_1$ 因子可以丢掉。]

\subsubsection{纳入磁物质}

我们还可以把带磁荷的 $(d-3)$ 维动力学激发
\begin{equation}
  M_i\,,\qquad 1\le i\le m\,, \tag{3.58}
\end{equation}
以有限张力加入到 阿贝尔规范理论中。设 $M_i$ 具有磁荷 $p_{i,a}\in\Z$，它们可以组织成一个 $m\times r$ 的荷矩阵 $P$，其元素为 $p_{i,a}$。这类带有磁激发的 阿贝尔规范理论，会作为描述超对称 CFT 以及更一般的超对称 QFT 的 Coulomb 相的低能有效理论出现。

这些磁激发会导致磁荷为
\begin{equation}
  p_{i,a}\,,\qquad 1\le a\le r\,. \tag{3.59}
\end{equation}
的 't Hooft 算符 $\cH_i$ 被完全屏蔽。未屏蔽 't Hooft 算符的群按下式计算：
\begin{equation}
  \mathsf{D}_{d-3}=\prod_{i=1}^{\min(m,r)}\Z_{p_i^D}\,, \tag{3.60}
\end{equation}
其中 $p_i^D$ 是 $P$ 的 Smith 标准形内部的元素
\begin{equation}
  \diag(p_1^D,p_2^D,\cdots,p_{\min(m,r)}^D)\subset \mathrm{SNF}(P)\,. \tag{3.61}
\end{equation}
$\mathsf{D}_{d-3}$ 的 Pontryagin 对偶就是磁 $(d-3)$-形式对称性群
\begin{equation}
  G^{(d-3)}=\widehat{\mathsf{D}}_{d-3}=\prod_{i=1}^{\min(m,r)}\widehat{\Z}_{p_i^D}\,. \tag{3.62}
\end{equation}

% [原文 p.37]
我们提醒读者注意我们的记号 (3.51)。

加入磁激发不影响电 1-形式对称性群 (3.50)。

\begin{examplebox}{例 3.7：把 $\SU(2)$ Higgs 化到 $\U(1)$}
再次考虑例 3.2。正如后面将要讨论的，UV 理论的磁 $(d-3)$-形式对称性群是平凡的：
\begin{equation}
  G^{(d-3)}=1\,. \tag{3.63}
\end{equation}
这一点在 IR 理论中可以通过下述观察重现：UV 的 $\SU(2)$ 理论的 't Hooft--Polyakov 磁单极子解，成为 IR 的 $\U(1)$ 理论的一个动力学磁激发 $M$。正如 (3.90) 式附近所讨论的，$\SU(2)$ 规范理论中的磁单极位形与 $\SU(2)$ 群流形中的圈相关联。与 't Hooft--Polyakov 磁单极子相关联的那个圈绕极大环面 $\U(1)$ 一圈，这意味着 $M$ 的 IR 磁荷为
\begin{equation}
  p=1\,, \tag{3.64}
\end{equation}
其结果是所有 IR 't Hooft 线算符都被完全屏蔽。
\end{examplebox}

\begin{examplebox}{例 3.8：把 $\SO(3)$ Higgs 化到 $\U(1)$}
现在把上一个例子中的 UV 规范群换成 $\SO(3)$。正如后面将要讨论的，UV 理论的磁 $(d-3)$-形式对称性群此时为
\begin{equation}
  G^{(d-3)}=\Z_2\,. \tag{3.65}
\end{equation}
为了在 IR 理论中重现这一点，注意 $M$ 的 IR 磁荷现在是
\begin{equation}
  p=2\,, \tag{3.66}
\end{equation}
因为绕极大环面 $\U(1)\subset \SU(2)$ 一圈的那个圈，会绕极大环面
\begin{equation}
  U(1)'=\frac{U(1)}{\Z_2}\subset SO(3) \tag{3.67}
\end{equation}
两圈。
\end{examplebox}

% [Tutor Note: p.37 例 3.7--3.8 对照的是“磁单极子绕圈数”：$\pi_1(\SU(2))=1$ 保证任何以 $\U(1)$ 为像的圈都可缩，于是 't Hooft--Polyakov 单极子能作为动力学激发存在；对 $\SU(2)$，绕 $\U(1)$ 极大环面一圈的圈在 $\SO(3)=\SU(2)/\Z_2$ 中要绕两圈才闭合，故 (3.64) 的 $p=1$ 变成 (3.66) 的 $p=2$。核对锚点：$\Z_p$ 中未屏蔽 't Hooft 线的步长就是 $p$，$p=1$ 给出全屏蔽（(3.63) 的平凡群），$p=2$ 给出 $\Z_2$（(3.65)）。]

\subsubsection{纳入双荷物质}

在时空维数 $d=4$ 时，上述图景会有一个推广，因为带电荷与带磁荷的算符具有相同的维数。考虑一个 $4d$ 的 阿贝尔规范理论，规范群为 $\U(1)^r$。此时 't Hooft 算符是线算符，它们与 Wilson 线算符组合起来，构成一般的双荷 \textbf{'t Hooft--Wilson 线算符}。

% [原文 p.38]
这些线算符可以被理论中的动力学粒子屏蔽，而这些粒子现在也可能带有双荷，即同时携带电荷与磁荷。令 $1\le i\le m$ 标记所有不同的粒子，其电荷为 $q_{i,a}$、磁荷为 $p_{i,a}$。这类带有双荷激发的 阿贝尔规范理论，会作为描述 $4d$ 超对称 CFT 以及更一般的 $4d$ 超对称 QFT 的 Coulomb 相的低能有效理论出现。

这些双荷激发导致电荷与磁荷分别为
\begin{align}
  q_{i,a}\,,&\qquad 1\le a\le r\,,\nonumber\\
  p_{i,a}\,,&\qquad 1\le a\le r\,, \tag{3.68}
\end{align}
的 't Hooft--Wilson 线算符 $\cD_i$ 被完全屏蔽。

利用动力学粒子的电荷与磁荷，我们定义一个 $m\times 2r$ 的荷矩阵 $C$，其第 $i$ 行是荷矢量
\begin{equation}
  C_i=(q_{i,1},q_{i,2},\cdots,q_{i,r},p_{i,1},p_{i,2},\cdots,p_{i,r})\,. \tag{3.69}
\end{equation}
我们然后计算
\begin{equation}
  \diag(c_1^D,c_2^D,\cdots,c_{\min(m,2r)}^D)\subset \mathrm{SNF}(C)\,, \tag{3.70}
\end{equation}
于是未屏蔽 't Hooft--Wilson 线算符的 阿贝尔群为
\begin{equation}
  \mathsf{D}_1=\prod_{i=1}^{\min(m,2r)}\Z_{c_i^D}\,. \tag{3.71}
\end{equation}
它的 Pontryagin 对偶是一个\textbf{双荷 1-形式对称性群}
\begin{equation}
  G^{(1)}=\widehat{\mathsf{D}}_1=\prod_{i=1}^{\min(m,2r)}\widehat{\Z}_{c_i^D}\,. \tag{3.72}
\end{equation}

% [Tutor Note: p.38 式 (3.68)--(3.72) 把电、磁两套荷合并进一个 $m\times2r$ 矩阵 $C$（(3.69) 的前 $r$ 列是电荷、后 $r$ 列是磁荷），因此 SNF 的上限从 $\min(m,r)$ 变成 $\min(m,2r)$——这是读者最容易漏看的一处。核对方法：当 $m=2r$ 且 $C$ 可逆时，$|\det C|$ 就是 $\Z^{2r}/\mathrm{im}(C)$ 的阶，可直接与 SNF 对角元之积 $\prod_i c_i^D$ 对比；对角元为 $0$ 的因子按 (3.51) 的约定给出 $\Z_0=\Z$，即未被屏蔽的连续族。]

\begin{examplebox}{例 3.9：Seiberg--Witten 模型中的双荷点}
考虑 $4d$ $\cN=2$ 纯 $\SU(2)$ 规范理论。我们后面将看到，该理论具有一个 1-形式对称性
\begin{equation}
  G^{(1)}=\Z_2\,. \tag{3.73}
\end{equation}
众所周知 [77]，该理论有一个有意思的真空 库仑分支，其中包含两个特殊点，即磁单极子点与双荷点。考虑双荷点处的 IR 理论，它是一个 $\cN=2$ SQED，具有 $\U(1)$ 规范群以及一个电荷为
\begin{equation}
  q_\phi=1 \tag{3.74}
\end{equation}
的物质场 $\phi$。此外我们还有一个有质量的 BPS 粒子，从
% [原文 p.39]
SQED 的观点看它是一个双荷，带电荷与磁荷
\begin{equation}
  q=1\,,\qquad p=2\,. \tag{3.75}
\end{equation}
把这些电荷收集起来，我们得到双荷荷矩阵
\begin{equation}
  C=\begin{pmatrix} 1 & 0\\ 1 & 2\end{pmatrix}\,, \tag{3.76}
\end{equation}
其 Smith 标准形为
\begin{equation}
  \mathrm{SNF}(C)=\begin{pmatrix} 1 & 0\\ 0 & 2\end{pmatrix}\,, \tag{3.77}
\end{equation}
这意味着 IR 的 1-形式对称性群同样是
\begin{equation}
  G^{(1)}=\Z_2\,. \tag{3.78}
\end{equation}
\end{examplebox}

% [Tutor Note: p.38--39 例 3.9 可直接核对：(3.76) 的 $2\times2$ 矩阵有 $\det C=1\cdot2-0\cdot1=2$，故 $\Z^2/\mathrm{im}(C)$ 的阶为 2；其 SNF 为 (3.77) 的 $\diag(1,2)$，代入 (3.71) 得 $\mathsf{D}_1=\Z_1\times\Z_2=\Z_2$，对偶后仍是 $\Z_2$，与 (3.73) 的 UV 结果匹配。这正是 Seiberg--Witten 理论在双荷点上“1-形式对称性沿 RG 流不变”的一个具体检验；作者所引 [77] 指原始的 Seiberg--Witten 解。]
